\cotefile{79}
\scanpage{235}
% Upper edge, fragment of the preceding strip
On trouve
\mspage{119 bis}

\unsure{graphes|graphe}, \inserted{avec} groupe d'homothéties $Z$ \unsure{qu'on se le donne|qu'on se donne|opérant librement}, \inserted{et satisfaisant la condition (41\unsure{*|''})}, \deleted{revient à celle} revient \deleted{à la} aux données \deleted{des choses} suivantes

\[
\left\{\begin{aligned}
&\text{a) un \deleted{gr.} }S^{\unsure{\sharp|\natural|\flat}},\ \text{une relation symétrique }R^\sharp\\
&\qquad\text{(i.e. un \uline{graphe} }\gamma^\sharp\text{) }\deleted{\unsure{\mathrm{sur}|\mathrm{entre}}}\ p_1,p_2:R^\sharp\rightrightarrows S^\sharp\\
&\text{b) }\deleted{\unsure{S}}\ \text{un }Z\text{-torseur relatif }S/S^\sharp\text{ sur }S^\sharp\\
&\text{c) un isomorphisme}\\
&\qquad p_1^*(S^{\deleted{\unsure{\sharp}}}/S^\sharp)\deleted{\unsure{\otimes}}\mathbin{\wedge}p_2^*(S/S^\sharp)\simeq\underset{\text{torseur trivial}}{\mathbb{1}_{Z_{S^\sharp}}}
\end{aligned}\right.
\tag*{(48.7)}
\]

% slanted
\margin{C'est \deleted{\unsure{essentiellement}} la relation symétrique sur $S$ associée à la structure de graphe}

Enfin, dans cette optique, la donnée d'un multigraphe $\pi$, \deleted{\unsure{avec|$\pi$}} avec groupe d'homothéties, revient aux données suivantes

\[
\left\{\begin{aligned}
&\text{a) Un graphe }\gamma^\sharp=(S^\sharp,\underbrace{R^\sharp\subset S^\sharp\times S^\sharp}_{\text{rel. sym.}})\\
&\text{b) Un }Z\text{-torseur relatif }S^{\deleted{\unsure{\sharp}}}/S^\sharp\\
&\text{c) Un }Z\text{-torseur absolu }\pi\\
&\text{d) un isomorphisme}\\
&\qquad p_1^*(S/S^\sharp)\wedge p_2^*(S/S^\sharp)\simeq\pi
\end{aligned}\right.
\tag*{(48.8)}
\]

Ainsi : la description (49.7), dont $i\in\pi$ définit alors une structure de graphes à homothéties sur $S$, comme ens de sommets, puisqu'il définit un iso $p_1^*(S/S^{\unsure{\sharp|\flat}})\wedge p_2^*(S/S^\sharp)\simeq\pi\xrightarrow[\text{via }i]{\sim}\mathbb{1}$.

\uline{Ex.} Si $M$ module libre de $\mathrm{rg}\,2$ sur $\Lambda$, $S=M^*/{\pm1}$, $\check{S}^{*}\simeq\check{M}^{*}/{\pm1}$ si $Z\deleted{\unsure{\mathrm{anneau}}}=\Lambda^*/{\pm1}$, \deleted{\unsure{on}} $\pi=(\det M)^{*}/{\pm1}$, on a $S^\sharp=M^*/\Lambda^*\subset\mathbb{P}(M)(\Lambda)$ (égalité si $\mathrm{Pic}(\Lambda)\simeq\unsure{0|0'}$), $R^\sharp\deleted{\unsure{=}}\subset S^\sharp\times S^\sharp$ est formé des \deleted{classes de couples} \deleted{\unsure{$(x,y)$}} $(x^\sharp,y^\sharp)$, avec $x,y\in M^*$, $x\wedge y\in(\det)^{*}$ --- i.e. les couples de sections du fibré projectif, distinctes en tt pt de la fibre. Le $S$-torseur relatif, sur \deleted{$\mathbb{P}(M)$} $S^\sharp$ correspond au fibré des \unsure{Serre-Hopf|semi-droites} modulo $\pm1$. L'iso. \deleted{\unsure{est}} d) est celui déduit de $\det M^*\simeq D^{\unsure{\wedge|1|*}}\mathbin{\wedge}_{\Lambda^*}D'^{\unsure{\wedge|1|*}}$ en passant à $\Lambda^*/{\pm1}$

% slanted
\margin{ou encore les couples de \textquotedblleft{}droites\textquotedblright{} $D,D'$ dans $M$ telles que $M=D\oplus D'$ --- d'où $\deleted{\unsure{\det}}\det M\simeq D\otimes D'$ donc $\det M^*\simeq D^\wedge\mathbin{\wedge}_{\Lambda^*}D'^\wedge$}

% archive sheet number: 234
